\chapter*{使用说明}
\pdfbookmark[chapter]{使用说明}{guide}

这份笔记按“设计模型如何变成可制造的产品”组织复习。前半部分讨论图形变换与几何表达，后半部分讨论实物数据如何回到模型，以及模型如何进入工艺设计、数控加工和3D打印。计算题应同时掌握公式、运算顺序与几何意义；简答题应能说明输入、处理过程、输出及适用条件。

\section*{材料范围与阅读顺序}
正文第1至7章对应马彦昭老师的七份课程讲义；第8、9章分别整理梁良老师的CAPP和APT补充讲义。页码均指原PDF从封面起计的物理页码。讲义中的章节编号有局部沿用旧编号的情况，本文按文件的第1至7章顺序编排。

\begin{enumerate}
\item \textbf{先建立主线：}读第1章，区分CAD、CAE、CAPP和CAM；用第3章理解几何、拓扑与设计意图的关系。
\item \textbf{再练计算：}学习第2章的齐次变换与裁剪，第4、5章的曲线曲面。每次先写向量约定、参数区间和边界条件，再代数值。
\item \textbf{最后串联应用：}学习第6至9章，比较逆向重建、增材制造和数控加工的数据、精度与工艺约束。
\item \textbf{用作业检验：}附录A保留讲义和独立作业图片中的原题；共同题合并，差异题分别标注。附录B提供跨章节速查和综合自测。
\end{enumerate}

\nt{本文的例题用于解释讲义方法；除附录A明确标注的原作业外，正文中的演算和情境题均为复习示例，不代表老师额外布置的作业或已确定的考试题。}

\section*{符号与精度约定}
图形变换统一用\textbf{列向量左乘}：$\widetilde{\boldsymbol p}'=T\widetilde{\boldsymbol p}$。若原讲义写成行向量右乘，应将矩阵转置后使用；连续操作先做的矩阵位于乘积右侧。旋转取右手正向，角度应明确旋转轴的正方向。

曲线以$\boldsymbol C(u)$表示，曲面以$\boldsymbol S(u,v)$表示。B样条\textbf{次数}记为$p$，\textbf{阶次}为$k=p+1$。原讲义采用$N_{i,k}$表示$k$阶基函数，本文以$N_{i,p}$表示$p$次基函数；两种下标不能直接混用。

讲义中的设备性能、企业案例和软件发展描述用于理解课程背景，不作为当前产品性能的承诺。正文优先保留原理和比较方法；对易混的公式从定义推导，并补全必要条件。配图从原PDF局部裁取，图注注明来源，编译不依赖原始下载目录。
